\documentclass[11pt]{article}
\usepackage[margin=0.88in]{geometry}
\usepackage[T1]{fontenc}
\usepackage{lmodern}
\usepackage{amsmath,amssymb,amsthm,mathtools}
\usepackage{booktabs,array,microtype,xcolor}
\usepackage[hidelinks]{hyperref}
\usepackage{fancyhdr}
\pagestyle{fancy}\fancyhf{}
\fancyhead[L]{\small Six-point homometry over fields}
\fancyhead[R]{\small Research note, 30 September 2026}
\fancyfoot[C]{\thepage}
\setlength{\headheight}{14pt}
\setlength{\parskip}{3pt}
\setlength{\parindent}{1.2em}
\newtheorem{theorem}{Theorem}
\newtheorem{lemma}{Lemma}
\newtheorem{corollary}{Corollary}
\renewcommand{\thetheorem}{\Alph{theorem}}
\renewcommand{\thecorollary}{\Alph{corollary}}
\setcounter{corollary}{3}
\newcommand{\F}{\mathbb F}
\newcommand{\Z}{\mathbb Z}
\newcommand{\TI}{\mathrm{T/I}}
\newcommand{\Om}{\Omega_K}
\newcommand{\status}[1]{\textsf{[#1]}}
\newcommand{\filepath}[1]{\path{#1}}
\title{Six-point homometry over fields:\\exact classification and counting of a classical family}
\author{}
\date{30 September 2026}
\hypersetup{pdftitle={Six-point homometry over fields: exact classification and counting of a classical family},pdfauthor={},pdfsubject={Direct algebraic classification of the Bloom image, with explicit computer-assisted cyclic total-count dependencies}}
\begin{document}
\maketitle\thispagestyle{empty}
\begin{abstract}
We classify the image of a classical two-parameter pair of homometric
six-point configurations under translation and inversion over fields of
characteristic different from $2,3,11$. Twelve explicit parameter symmetries
account for every identification, including coincidences after projection.
Centered moments reduce the classification to a cubic root-permutation
problem. Over finite fields this gives exact counts; a single endpoint
determines its partner except in characteristic $31$, where an explicit
counterexample is given. These are direct algebraic results. Total counts
for six-subsets of prime cyclic groups additionally use an earlier
computer-assisted completeness theorem and finite census. No historical
novelty, external peer review, or proof-assistant certification is asserted.
\end{abstract}

\section{The problem and the counted object}
For a finite subset $A$ of the additive group of a field $K$, define its
directed autocorrelation by
\[
 c_A(h)=\#\{(x,y)\in A^2:x-y=h\}\qquad(h\in K).
\]
Sets are \emph{homometric} if their autocorrelations agree. This is a
reconstruction problem from a difference multiset; it is relevant to
crystallographic phase retrieval, sparse support recovery, and the cyclic
beltway problem, as well as to musical interval content. For $K=\F_p$ with
$p$ an odd prime, it is equivalent to equality of cyclic interval-class
vectors. For $K=\F_{p^r}$ with $r>1$, the additive group is
$(\Z/p\Z)^r$, not the cyclic group $\Z/p^r\Z$.

Two sets are \emph{rigidly equivalent} if $B=\epsilon A+t$, where
$\epsilon\in\{1,-1\}$ and $t\in K$; write $[A]$ for a rigid class.
General nonzero scalar multiplication is \emph{not} identified. With
$v=(a,b)^T\in K^2$, consider the classical configurations
\begin{align*}
 X_v&=\{0,a,b-2a,2b-2a,2b,3b-a\},\\
 Y_v&=\{0,a,b+2a,2b-a,2b+a,3b-a\}.
\end{align*}
Let $\Om$ consist of parameters for which each list has six distinct
entries, and put $\Phi(v)=\{[X_v],[Y_v]\}$. The endpoints of this
unordered pair may be translated and inverted independently.
We call the image of $\Phi$ the \emph{Bloom image}; this is terminology
for a classical family, not a priority assertion.

Postpischil--Gilbert give these formulas and their triangular projection
geometry \cite[\S\S3,7]{pg}. The issue here is the exact quotient after
projection into a field. Formal symmetries alone do not exclude accidental
identifications. The moment argument below supplies that missing step,
without a distinct-distance assumption or a real ordering. Theorems A--C
concern this specified image. They do not by themselves classify every
homometric six-subset of every additive group.

\newpage
\section{Direct algebraic results}
Throughout Theorems A--C, $K$ has characteristic different from $2,3,11$.
The results are \status{PROVED}, in-house, with a separate adversarial
review recorded in the accompanying project note. The proofs below are
algebraic; finite computations are supporting checks.

\begin{theorem}[exact image classification]\label{thm:A}
Every $v\in\Om$ gives two homometric, different rigid classes. Put
\[
 R=\begin{pmatrix}0&-1\\1&-1\end{pmatrix},\qquad
 T=\begin{pmatrix}0&1\\1&0\end{pmatrix},\qquad
 G=\{\pm I,\pm R,\pm R^2,\pm T,\pm RT,\pm R^2T\}.
\]
Then $G$ is a group of order twelve, isomorphic to $S_3\times C_2$,
which acts freely on $\Om$ and preserves $\Phi$. For $v,w\in\Om$,
\[
 \Phi(v)=\Phi(w)\quad\Longleftrightarrow\quad w\in Gv
 \quad\Longleftrightarrow\quad
 \bigl(s(v),e(v)^2\bigr)=\bigl(s(w),e(w)^2\bigr),
\]
where $s=a^2-ab+b^2$ and $e=ab(a-b)$.
Thus every unordered pair has exactly twelve parameter representatives.
\end{theorem}

\begin{theorem}[finite-field enumeration]\label{thm:B}
For $K=\F_q$ in the stated characteristic, the number of distinct
unordered pairs in the Bloom image is
\[
 B(q)=\begin{cases}
 (q-1)(q-5)/12,&\operatorname{char}K=5,\\[2pt]
 (q-1)(q-7)/12,&\operatorname{char}K=7,\\[2pt]
 (q-1)(q-11)/12,&\operatorname{char}K\ge13.
 \end{cases}
\]
In particular $B(p)=(p-1)(p-11)/12$ for every prime $p\ge13$.
The displayed configurations give no six-element pair in
$\F_2,\F_3,\F_5,\F_7,\F_{11}$.
\end{theorem}

The last assertion is a support observation, not an extension of the
classification proof to characteristic $11$. Extension fields of
characteristics $2,3,11$ are outside Theorems A--B.

\begin{theorem}[unique partner within the image]\label{thm:C}
If $\operatorname{char}K\ne31$ as well, two distinct pairs in the Bloom
image cannot share a rigid class. Consequently the graph with rigid
classes as vertices and Bloom pairs as edges is a disjoint union of edges.
\end{theorem}

Theorem C permits additional homometric partners outside the Bloom image.
Characteristic $31$ is a genuine exception: Section~\ref{sec:partner}
gives two different Bloom pairs sharing one endpoint. Theorems A--B
continue to hold there.

For example, $B(13)=2$, $B(31)=50$, and $B(1009)=83\,832$.
The counts in the noncyclic additive fields $\F_{25},\F_{49},\F_{169}$
are respectively $40,168,2212$. These field counts retain translation
and inversion as their sole equivalences. A quotient by all field
scalars or by all additive-group automorphisms would be a different problem.

\newpage
\section{Differences, collisions, and the twelve symmetries}
The fifteen differences of either configuration, taken up to sign and
with multiplicity, are the following common list:
\begin{gather*}
 a,\ b-2a,\ 2b-2a,\ 2b,\ 3b-a,\ b-3a,\ 2b-3a,\ 2b-a,\\
 3b-2a,\ b,\ b+2a,\ 2b+a,\ 2a,\ b-a,\ b+a.
\end{gather*}
Replace each entry $\ell$ by $\ell,-\ell$ and include the six diagonal
zero differences. This gives equal directed autocorrelations whenever
both supports have size six. It is an integer coefficient identity;
repeated nonzero differences cause no difficulty. Reduction into an
additive group preserves the identity for the labeled lists, but support
collisions must first be excluded to interpret them as six-element sets.

As $2$ is invertible, collisions occur exactly on the union of the
following twelve kernels:
\begin{gather}
 a,\ b,\ b-a,\ b+a,\ b-2a,\ b+2a,\ b-3a,\notag\\
 2b-a,\ 2b+a,\ 2b-3a,\ 3b-a,\ 3b-2a.\label{eq:lines}
\end{gather}
Some kernels coincide in characteristics $5$ and $7$; the union remains
exact. In particular, for every admissible parameter,
\[
 e=ab(a-b)\ne0,\qquad d=(a+b)(2a-b)(a-2b)\ne0.
\]

The relations $R^3=T^2=I$, $TRT=R^{-1}$, and central $-I$ show
that $G\cong S_3\times C_2$. Its twelve matrices are distinct in
characteristic different from $2,3$. Direct substitution gives the
following identities of sets:
\begin{center}
\begin{tabular}{@{}lll@{}}\toprule
$M$ & $X_{Mv}$ & $Y_{Mv}$\\\midrule
$I$ & $X_v$ & $Y_v$\\
$R$ & $X_v+2a-2b$ & $Y_v+a-3b$\\
$R^2$ & $X_v-2b$ & $Y_v-2a-b$\\
$T$ & $Y_v+a-2b$ & $X_v+2a-b$\\
$RT$ & $Y_v-a$ & $X_v-a$\\
$R^2T$ & $Y_v-a-2b$ & $X_v+a-3b$\\\bottomrule
\end{tabular}
\end{center}
For $-M$, negate the sets and translations on the right. Hence the
whole group preserves admissibility and $\Phi$.

To prove freeness, note that $\det(R-I)=\det(R^2-I)=3$ and
$\det(R+I)=\det(R^2+I)=1$. Thus $\pm R,\pm R^2$, and $-I$
have only the zero fixed vector. The fixed lines of
$T,-T,RT,-RT,R^2T,-R^2T$ are, respectively,
\[
 b=a,\quad b=-a,\quad a=0,\quad a=2b,\quad b=0,\quad b=2a.
\]
Each is excluded by \eqref{eq:lines}; zero is excluded as well.
Every orbit therefore has twelve parameters. This still leaves open
whether distinct orbits can become the same pair after projection.

\newpage
\section{Integral centered moments recover the orbit invariants}\label{sec:moments}
For a six-element set $A\subset K$, set
\[
 \mu_A=\frac16\sum_{x\in A}x,\qquad
 m_j(A)=\sum_{x\in A}\bigl(3(x-\mu_A)\bigr)^j.
\]
Translation preserves each $m_j$ and inversion multiplies it by
$(-1)^j$. Thus $m_2,m_3^2,m_6$ are invariants of a rigid class;
their endpoint sums are invariants of an unordered pair.

All identities needed here can be checked over $\Z$. Indeed
$3(x-\mu_A)=3x-\frac12\sum A$, while
$\sum X_v=-4a+8b$ and $\sum Y_v=2a+8b$. The scaled centered
coordinates are the integral linear forms
\begin{align*}
 X:\ &(2a-4b,\ 5a-4b,\ -4a-b,\ -4a+2b,\ 2a+2b,\ -a+5b),\\
 Y:\ &(-a-4b,\ 2a-4b,\ 5a-b,\ -4a+2b,\ 2a+2b,\ -4a+5b).
\end{align*}
Raising these six forms to the stated powers and summing gives
\begin{align}
 d^2+27e^2&=4s^3,\label{eq:disc}\\
 m_2(X_v)=m_2(Y_v)&=66s,\label{eq:m2}\\
 m_3(X_v)&=6d-243e,&m_3(Y_v)&=6d+243e,\label{eq:m3}\\
 m_6(X_v)&=23946s^3+6399e^2-4860de,\label{eq:m6x}\\
 m_6(Y_v)&=23946s^3+6399e^2+4860de.\label{eq:m6y}
\end{align}
For an explicit expansion check, the coefficients of the summed sixth
moments in the monomial order $a^6,a^5b,\ldots,b^6$ are
\[
 (47892,-143676,300150,-360840,300150,-143676,47892),
\]
the same as those of $47892s^3+12798e^2$.

Define invariant endpoint sums
\begin{align*}
 C_3&=m_3(X_v)^2+m_3(Y_v)^2-288s^3=116154e^2,\\
 C_6&=m_6(X_v)+m_6(Y_v)-47892s^3=12798e^2.
\end{align*}
The integer relations
\[
 116154=162\cdot717,\quad12798=162\cdot79,\quad
 118\cdot79-13\cdot717=1
\]
yield the exact elimination identity
\begin{equation}
 \boxed{118C_6-13C_3=162e^2.}\label{eq:eliminate}
\end{equation}
If $\Phi(v)=\Phi(w)$, equation \eqref{eq:m2} first recovers the common
$s$, since $66$ is invertible in the stated characteristic. The endpoint
sums then recover $C_3,C_6$, and \eqref{eq:eliminate} recovers $e^2$,
since $162$ is invertible. Using both moments avoids exclusions at
characteristics $79$ and $239$, where one individual coefficient vanishes.
There is no division by $s$: isotropic parameters $s=0$ are included.

\newpage
\section{Cubic root permutations and finite-field counting}
Suppose $s(v)=s(w)$ and $e(v)^2=e(w)^2$. As the characteristic is odd,
$e(w)=e(v)$ or $e(w)=-e(v)$. In the latter case replace $w$ by $-w$,
an element of $G$, so that the values of $e$ agree. The identity
\begin{equation}
 (z-a)(z+b)(z-b+a)=z^3-sz-e\label{eq:cubic}
\end{equation}
then gives the same multiset of roots for the two ordered triples
$(a,-b,b-a)$. This includes roots with multiplicity: uniqueness of
factorization over a field is sufficient.

The map $R$ cyclically permutes this ordered triple; $-T$ transposes
its first two entries and fixes the third. They generate all six
permutations. The first two entries recover $a,b$, so the aligned $w$
is in the group generated by $R,-T$. Restoring any negation places the
original $w$ in $Gv$. Conversely the explicit symmetries preserve $\Phi$
and its recovered invariants. Thus equality of the invariants, equality
of pairs, and membership in one $G$-orbit are equivalent.

The two endpoints cannot collapse to one rigid class. From \eqref{eq:m3},
\[
 m_3(X_v)^2-m_3(Y_v)^2=-5832de\ne0.
\]
The factors $d,e$ are nonzero on $\Om$, and $5832$ has only $2,3$
as prime factors. Rigid equivalence would preserve the squared third
moment, a contradiction. This proves Theorem A in full, including the
absence of accidental projected identifications.

For Theorem B, the primitive normals of the collision lines are
\begin{gather*}
 (1,0),(0,1),(-1,1),(1,1),(-2,1),(2,1),(-3,1),\\
 (-1,2),(1,2),(-3,2),(-1,3),(-2,3).
\end{gather*}
The nonzero absolute pairwise determinants are exactly
$\{1,2,3,4,5,7,8\}$. Thus the twelve lines are distinct in
characteristic at least $13$. Their projective slopes $b/a$ are
\[
 \infty,0,1,-1,2,-2,3,1/2,-1/2,3/2,1/3,2/3.
\]
In characteristic $5$ these reduce to all six directions of
$\mathbb P^1(\F_5)$; in characteristic $7$, to all eight directions of
$\mathbb P^1(\F_7)$. They remain six or eight distinct lines over
an extension field, since distinct base-field slopes remain distinct.

Let $d_p$ be this number of lines over $K=\F_q$. Each contains $q-1$
nonzero vectors, and distinct lines meet only at zero. Therefore
\[
 |\Om|=q^2-1-d_p(q-1)=(q-1)(q+1-d_p).
\]
With $d_5=6$, $d_7=8$, and $d_p=12$ for $p\ge13$, divide by twelve
using Theorem A to obtain the three formulas in Theorem B. For
$\F_{11}$ the slopes cover all twelve projective directions, so there
is no admissible parameter. The fields $\F_2,\F_3,\F_5$ have fewer
than six elements; $\F_7$ is already covered by the direction argument.

\newpage
\section{A single endpoint and the characteristic-31 obstruction}\label{sec:partner}
A rigid class occurring at either endpoint first recovers $s$ from $m_2$.
For an endpoint $A=X_v$, equations \eqref{eq:disc}--\eqref{eq:m6y} give
\begin{equation}
 \begin{pmatrix}m_3(A)^2-144s^3\\m_6(A)-23946s^3\end{pmatrix}
 =\begin{pmatrix}58077&-2916\\6399&-4860\end{pmatrix}
 \begin{pmatrix}e^2\\de\end{pmatrix}.\label{eq:endpoint}
\end{equation}
For $A=Y_v$ the same matrix applies with second unknown $-de$.
Its determinant is
\[
 58077(-4860)-(-2916)6399
 =-263594736=-2^4\,3^{12}\,31.
\]
Outside characteristics $2,3,11,31$, the endpoint therefore determines
$s,e^2$, regardless of which side it occupies. Two Bloom pairs sharing
an endpoint have the same invariants and hence, by Theorem A, are the
same pair. This proves Theorem C. It again requires no division by $s$
or any data-dependent value.

The exceptional characteristic is necessary. In $\F_{31}$ take
$v=(2,12)$ and $w=(8,17)$, and set
\begin{align*}
 A&=\{0,1,3,8,12,18\},\\
 B&=\{0,1,3,10,14,26\},\\
 C&=\{0,1,4,10,12,17\}.
\end{align*}
The raw configurations are
\begin{align*}
 X_v&=\{0,2,3,8,20,24\},&Y_v&=\{0,2,3,16,22,26\},\\
 X_w&=A,&Y_w&=\{0,2,8,11,12,26\}.
\end{align*}
Their exact transports are
\[
 A=3-Y_v=X_w,\qquad B=3-X_v,\qquad C=12-Y_w.
\]
Both parameters are admissible, so $\{[A],[B]\}$ and
$\{[A],[C]\}$ are Bloom pairs. All three sets have each cyclic
interval $1,\ldots,15$ exactly once. Their sorted multisets of
successive circular gaps are, respectively,
\[
 (1,2,4,5,6,13),\qquad (1,2,4,5,7,12),\qquad (1,2,3,5,6,14).
\]
Translation rotates a circular gap sequence and inversion reverses it;
the different gap multisets exclude rigid equivalence. Thus the two
pairs genuinely differ and share exactly the stated class.

Moreover, $s(v)=s(w)=0$ but $e(v)^2=2$ and $e(w)^2=8$ in $\F_{31}$.
This exhibits failure of endpoint recovery, not merely a singular matrix
in one method of proof. Pair classification and the count $B(31)=50$
remain valid. The existence of this three-class subfamily alone does not
establish the complete size-five family asserted by the finite census
in the next section.

\newpage
\section{Total counts in prime cyclic groups: inherited dependencies}\label{sec:totals}
The direct algebra above counts the Bloom image. To identify it with
\emph{all} nontrivial six-point homometry in $\Z/p\Z$ at large primes,
we use a separate completeness input from the Homometry Program:

\medskip\noindent\textbf{Input G and its large-prime corollary.}
The earlier general six-point generation theorem is \status{PROVED},
in-house, \emph{computer-assisted}. Its large-prime corollary states
that every cyclic six-point pair of different rigid classes is an
admissible modular Bloom pair when every prime divisor of the modulus
exceeds $131$. This uses the exact line-classification solver obligation
and the bounded universal-presentation/torsion certificate. The torsion
order is at most $135$; the least prime above $131$ is $137$, so the
torsion image vanishes. Integer lifts then reduce the free image to the
line classification. Those computational obligations are inherited,
not proved or replaced by the present moment argument.

\begin{corollary}[large-prime totals, assuming Input G]
For every prime $p>131$, the number of unordered nontrivial six-point
pair edges, and the number of maximal nontrivial six-point families in
$\Z/p\Z$, are both $(p-1)(p-11)/12$. Each family has exactly two
rigid classes.
\end{corollary}
\begin{proof}
Input G identifies every pair with the Bloom image. Theorem B counts
its edges, and Theorem C makes the graph a disjoint union of edges,
since these primes differ from $31$.
\end{proof}

\begin{corollary}[all prime totals, with finite census input]
In addition to Input G, assume the earlier complete two-method six-set
census through modulus $135$. For primes $p\ge13$, put
$B(p)=(p-1)(p-11)/12$ and let $[p=r]$ be $1$ if $p=r$ and $0$
otherwise. The total numbers of pair edges and maximal nontrivial
families are
\begin{align*}
 P_6(p)&=B(p)+8[p=17]+9[p=19]+11[p=23]+20[p=31],\\
 F_6(p)&=B(p)+8[p=17]+9[p=19]+11[p=23]+11[p=31].
\end{align*}
Every family has size two except one size-five family in $\Z/31\Z$.
Primes below $13$ have no nontrivial six-set family.
\end{corollary}
\begin{proof}
Corollary D covers primes above $131$. For the remaining primes the
complete census gives extra edge counts $8,9,11,20$ at $17,19,23,31$,
and zero elsewhere. At $31$ there are sixty size-two families and one
size-five family, hence $61$ families and $60+\binom52=70$ edges;
subtracting $B(31)=50$ gives the two different corrections. At
$17,19,23$ all families have size two. Primes below $7$ support no six-subsets,
and the complete checks at $7,11$ give no nontrivial families.
\end{proof}

Corollaries D--E are \status{PROVED}, in-house, \emph{computer-assisted}
under these recorded inputs. The small-prime replay checks canonical
classes, autocorrelations, disjoint membership, family sizes and source
hashes; its completeness is inherited from the original enumeration.
It is not a new exhaustive census or a pure algebraic classification
of the small-prime additional pairs.

\newpage
\section{Verification, provenance, and remaining scope}\label{sec:provenance}
The mathematical source is
\filepath{notes/2026-09-30-six-prime-theorem.md}; its separate attack is
\filepath{notes/2026-09-30-six-prime-review.md}.
An auxiliary coefficient-symmetry derivation and an independent polynomial
kernel check all displayed identities, collision determinants, and the
single-endpoint determinant. The proofs in Theorems A--C require only
those written identities and elementary field algebra.

Two independent implementations provide \status{COMPUTED} controls.
The direct point enumerator canonicalizes each endpoint under translation
and inversion and retains every parameter fiber, without using the formula
or symmetry group. The invariant implementation partitions parameters by
$(s,e^2)$ without constructing point classes. The \emph{entire partitions}
agree at every prime $13\le p\le251$ and at $1009$, and at the additive
fields $\F_{25},\F_{49},\F_{169}$. These checks support the algebraic
proof and do not supply general completeness. Tests cover isotropic
parameters, the coefficient exceptions $79,239$, and the explicit
characteristic-$31$ shared endpoint.

The retained complete six-set census has its original two independent
methods. Replaying its small-prime payloads validates the corrections in
Corollary E while retaining the original completeness dependency. Full
source and evidence hashes, native compiler diagnostics, exact export
commands and page inspection are saved alongside this PDF.
Useful reproduction commands in the pinned project environment are:
\begin{quote}\small\ttfamily
.venv/bin/python tests/run\_tests.py\\
.venv/bin/python tests/test\_six\_prime\_bloom.py\\
.venv/bin/python tests/test\_six\_bloom\_invariants.py\\
.venv/bin/python src/six\_bloom\_invariants.py
\end{quote}
The large-prime input is in the earlier working paper's section
``Integer shadows and the large-prime boundary'';
\filepath{notes/2026-09-30-six-paper.tex}. Its proof and certificate
obligations are not replaced by this standalone paper.

\paragraph{Prior art and limits.}
The family and its triangular geometry are classical \cite{pg}.
Ranieri et al. display the same six-point matrices in the collision-free
real-line setting \cite[\S IV.A]{ranieri}; that scope is distinct from
modular coincidences. The earlier small-prime family totals
$2,16,21,33$ at $13,17,19,23$ appear in Jedrzejewski--Johnson
\cite[\S3]{jj}; family counts and pair-edge counts differ when a family
has more than two members. The focused primary-source literature log is
\filepath{notes/2026-09-30-six-prime-literature.md}. Database access
failures and unread historical sources remain recorded; no novelty is
claimed. Composite cyclic moduli, a complete classification over noncyclic
additive fields, and extension fields of characteristics $2,3,11$ remain
outside this paper. External specialist review is still needed before
public dissemination; no submission or external publication has occurred.

\begin{thebibliography}{9}\small
\bibitem{pg} E. Postpischil and P. Gilbert,
There Are No New Homometric Golomb Ruler Pairs with 12 Marks or Less,
\emph{Experimental Mathematics} \textbf{3}(2) (1994), 147--152.
\href{https://doi.org/10.1080/10586458.1994.10504286}{doi:10.1080/10586458.1994.10504286}.
Primary sections used: 3 and 7.
\bibitem{ranieri} J. Ranieri, A. Chebira, Y. M. Lu and M. Vetterli,
\emph{Phase Retrieval for Sparse Signals: Uniqueness Conditions},
\href{https://arxiv.org/abs/1308.3058v2}{arXiv:1308.3058v2} (2013), \S IV.A.
\bibitem{jj} F. Jedrzejewski and T. Johnson,
\emph{The Structure of Z-Related Sets},
\href{https://arxiv.org/abs/1304.6608v1}{arXiv:1304.6608v1} (2013), \S3;
MCM 2013, LNCS 7937, 128--137.
\end{thebibliography}
\end{document}
